\section{Source Reconciliation, Record Boundaries and Transfer Objects}\label{app:source-reconciliation}
This targeted reporting audit revisits all 11 formerly unverified family R4/R5 cells and their affected record cells. It retrieves the official versioned HTML of six papers and records content hashes, section/table locators, old/new status and a rationale. Three further versioned papers support the positive reuse and holdout examples. It adds no model calls, customers or behavioral experiments. The remaining questions concern what the published evaluations bind; a new Sol run cannot establish the training provenance of somebody else's reported result.

\subsection{A bounded reading rule}
B requires the requested connection to be explicit. P applies when relevant training/profiling and defended results are described but their condition-specific connection is incomplete. U means the requested connection is not reported in the inspected methods/result passages; it does not mean the authors never performed it. A remains inaccessible source, and ? remains a binding not adjudicated within the reading scope. For R5, an undefended-versus-defended score or two detector architectures do not establish clean-trained versus protected-trained attacks on identical defended tests. We do not infer B from an implementation alone without its result binding.

\begin{longtable}{@{}p{.18\textwidth}p{.10\textwidth}p{.19\textwidth}p{.45\textwidth}@{}}
\caption{All 11 family-cell rereads. Each old code is ?. Source versions are pinned in the receipt file.}\label{tab:r26-source-review}\\
\toprule Family & New R4/R5 & Inspected locator & Evidence and remaining link\\\midrule\endfirsthead
\toprule Family & New R4/R5 & Inspected locator & Evidence and remaining link\\\midrule\endhead
\bottomrule\endfoot
Whisper Leak~\cite{whisperleak2025} & P/U & 3.2; 6.1-6.3; Tables 3-5; Figure 7 & Training and prompt splits are specified globally; padding uses new captures and batching/injection transform traces. These passages do not map protected training partitions to each mitigation result, or give clean/protected fits on identical defended tests.\\
Time Will Tell~\cite{timewilltell2024} & P/U & V-B; VII-B & The attack fits an input-length-dependent token-count threshold. The fixed-word instruction is evaluated, but its threshold reprofiling and a paired clean/protected-fit contrast are not specified in these passages.\\
Speculative decoding~\cite{wei2024speculation} & P/U & 4.3-4.4; 6.1; Appendix D & Fingerprint profiling and padded/aggregated attack outcomes are described. The mitigation result does not identify whether reference profiles are rebuilt on protected traces, nor compare both fits on a shared defended test population.\\
Cloud traffic benchmark~\cite{pouryousef2026traffic} & U/U & III-C; IV; V-C1; Figure 9(a) & The mitigation discussion evaluates fingerprints learned at a fixed configuration under rewritten prompts and changed temperatures. Protected-condition refitting and a matched clean/protected-fit contrast are not reported in this scope.\\
SparLeak~\cite{chen2026sparleak} & P/U & 6; 7.1; 7.5; Table 4 & Profiling-based attacks and replacement-intensity ASR/ROUGE outcomes are specified. The mitigation table does not bind the attack-training distribution to each intensity or report matched clean/protected-fit tests.\\
AHB~\cite{zhu2026turing} & P (retained)/U & 4.2.2; 5.1.2; Table 1; E.1-E.3 & Different detectors on a humanization row are a detector comparison, not a clean/protected-fit contrast. The inspected paper scope and archived static call-chain review do not report that matched contrast; condition-specific code fitting remains partial run evidence.\\
\end{longtable}


In the archived source-scope reconciliation, four family cells became P and seven U; no additional B was inferred. S36 documents the subsequent record-derived family projection, which preserves the six unresolved record bindings rather than inheriting these broader judgments. Nineteen record cells are updated consistently, including AHB's offline fitting as partial under the same code-versus-run distinction. Six other record-level ? cells remain in NetEcho, Early Bird and Dangol scopes, distinct from the 11 family cells reviewed here. Current record cells appear in Section~\ref{app:record-joins}; the artifact records the 19 exact changes. These targeted decisions resolve the reporting category, not the source's unpublished execution history.

\subsection{Effect on the quantitative synthesis}
All supported joins remain unchanged: inference--utility 4/22, inference--cost 5/22, matched adaptation--inference 2/17, triple join 2/20, and all five bindings 0/14. Query masks still start from the same 33 mitigation records. The triple package remains CachePrune's two outcomes at one shared configuration. The pair census remains 741 pairs, 133 related, with 17 descriptive matches, 105 mismatches and 11 unresolved pair contracts. Family-cell reporting and pair-contract status are separate objects.
\begin{table}[ht]\centering\small
\caption{Current query stress counts after source rereading. Only B is supported; weaker columns remain counterfactual. S33 preserves the earlier snapshot.}\label{tab:r26-current-stress}
\begin{tabular}{@{}lrrrrr@{}}\toprule Query & Eligible & B & B/P & B/? & B/P/?\\\midrule
Inference + utility & 22 & 4 & 6 & 4 & 6\\
Inference + cost & 22 & 5 & 11 & 5 & 11\\
Matched adaptation + inference & 17 & 2 & 2 & 6 & 6\\
Inference + utility + cost & 20 & 2 & 2 & 2 & 2\\
All five bindings & 14 & 0 & 0 & 0 & 0\\
\bottomrule\end{tabular}\end{table}


\subsection{Recordization and extraction coverage}
We preserve a hierarchy: family/version $\to$ comparison record $\to$ model/workload/parameter coordinates $\to$ measured outcomes. A comparison record splits when the primitive, execution mode, receiver/target contract or measured cost object changes. Multiple numerical cells within a common record are not additional independent observations. A join must match its configuration coordinates; two endpoints can share a quality/cost package without becoming a paired effect. The extraction frame is the 20 mitigation families' scoped evaluations and linked methods/outcomes plus six contextual controls. It is not an enumeration of all numeric cells or all attack experiments in the 82-family map.

\begin{itemize}
\item Early Bird: split the prefix-recovery oracle from semantic-cache processor cost (V-A versus V-B).
\item Dangol: split transformed captures from a live OHTTP loopback (9.8--9.9 versus 9.7); a common workflow target permits description, not a common execution denominator.
\item SparLeak: retain model/intensity coordinates within the random-replacement record; Table 4's ASR and ROUGE stay paired by row.
\item CachePrune: split direct and contextual recovery targets, but deduplicate the shared default quality/cost package.
\item OTRO: keep measured unpadded execution separate from a padded operating point and its derived cost.
\end{itemize}
These rules articulate the existing ledger retrospectively; they are not a claim of preregistered extraction. The released record-scope file gives every one of the 39 units, its source locator, retained parameter coordinates, configuration package and question cohort.

\begin{longtable}{@{}p{.19\textwidth}p{.08\textwidth}p{.08\textwidth}p{.56\textwidth}@{}}
\caption{Declared question groups. Full record identities are in the released scope file.}\label{tab:r26-cohorts}\\
\toprule Question group & Records & Pairs & Membership rule\\\midrule\endfirsthead
\toprule Question group & Records & Pairs & Membership rule\\\midrule\endhead
\bottomrule\endfoot
traffic envelope & 12 & 66 & Network/API destination, size, timing, count or delivery shaping; includes WRA destination cover, excludes cache probes and content-only identity.\\
cache service & 10 & 45 & Cache hit/reuse probes, cache policy and its component costs; excludes local memory observation without a service-cache probe.\\
local memory & 6 & 15 & Memory access/loading/sparsity/obliviousness and associated component costs; excludes network metadata and service-cache timing probes.\\
behavioral humanization & 3 & 3 & Site/mobile behavioral-trajectory edits or executed GUI humanization; excludes network-destination cover such as WRA.\\
generator fingerprint & 3 & 3 & Prompt/decoding or route changes with model-identity outcomes; excludes topic/content recovery.\\
fixed instruction & 2 & 1 & Fixed instruction/guardrail contrasts; source target and observer still require alignment.\\
access control & 1 & 0 & FP-Agent contextual gate; singleton, no related pair.\\
contention noise & 1 & 0 & Behind Bars contextual channel; singleton, no related pair.\\
filter created view & 1 & 0 & ML privacy channels contextual filter; singleton, no related pair.\\
\end{longtable}


Question assignment uses the primary extracted intervention/receiver surface, then the intended comparison. Network destination cover belongs to traffic-envelope even when it is implemented with a decoy agent; site/mobile behavioral edits belong to humanization. Fixed instructions and generator-identity countermeasures retain their declared narrower questions. Cost-only records stay with the mechanism they would otherwise be used to price. A record that could answer another question needs an explicit alternate-cohort audit; outside-cohort status is not evidence of incompatibility. The six nonsingleton cohorts generate all 133 related pairs; three singleton contextual groups generate none. The full cross-product remains available for alternative queries.

\subsection{What the holdout index evaluates}
The 19 H entries across six axes locate reported separations in the archived extraction. A family-level H does not itself bind a defense to that split. Table~\ref{tab:r26-h-objects} lists the entries and whether the family also has a mitigation record. The three model entries concern attack/identification generalization outside the 20 mitigation families. The two topology entries likewise concern attacker generalization: Dangol's new-composition recognition and the traffic benchmark's graph/star classifier transfer. Neither is a demonstrated source-calibration transfer of the defense. The other 14 entries retain their source-qualified split descriptions; classifying their defended-condition linkage requires the relevant record, rather than a blanket attack-only or defense-transfer label.

Defense-conditioned variation is still useful. Time Will Tell's shared fixed-length instruction has model-dependent outcomes; Whisper Leak's mitigation results remain provider/condition-specific. In these cases retain the tested rule and residual attack outcome, rather than relabel multiple conditions as a held-out defense-calibration experiment. For actual calibration transfer, specify source training, frozen parameters, target testing, local recalibration and attacker fitting separately. The artifact's H-entry ledger carries each original locator and qualification; the 82-family discovery matrix remains unchanged.
\begin{longtable}{@{}p{.25\textwidth}p{.16\textwidth}p{.10\textwidth}p{.40\textwidth}@{}}
\caption{Source-qualified H entries. The three model and two topology splits evaluate attacker generalization; other entries retain the indexed split description. The H flag alone does not establish a defended holdout or calibration transfer. M indicates mitigation-audit membership.}\label{tab:r26-h-objects}\\
\toprule Family & Axis & M & Archived qualification\\\midrule\endfirsthead
\toprule Family & Axis & M & Archived qualification\\\midrule\endhead
\bottomrule\endfoot
WRA~\cite{jeong2026wra} & Task/input/user & yes & Prompt tests and synthetic persona split; version reconciliation open\\
FP-Agent~\cite{wang2026fpagent} & Task/input/user & no & Separate task holdout; visitor/host separation unresolved\\
Prompt length~\cite{weiss2024prompt} & Task/input/user & no & OOD inputs; browser/API and day/night buffering controls\\
InputSnatch~\cite{inputsnatch2024} & Task/input/user & no & Held-out and generated target prompts; candidate-prior dependence\\
Whisper Leak~\cite{whisperleak2025} & Task/input/user & yes & Unique target-prompt holdout; model/provider comparisons\\
Prompt Echo~\cite{prompt_echo2026} & Task/input/user & no & Overlap-filtered prompts; models evaluated separately\\
I Know What You Said~\cite{gao2025iknow} & Task/input/user & yes & Prompt overlap filtering; model/hardware/framework comparisons\\
Speculation~\cite{wei2024speculation} & Task/input/user & yes & New traces and OOD paraphrases; remote/load conditions\\
Workflow traffic~\cite{dangol2026workflow} & Topology & yes & Leave-composition-out; stored and live populations separated\\
Traffic benchmark~\cite{pouryousef2026traffic} & Topology & yes & Topology shift distinct from feature-view restriction\\
ProcGrep~\cite{oderinwale2026procgrep} & Task/input/user & no & Task-grouped attribution under scaffold regimes\\
LLMs Have Rhythm~\cite{alhazbi2025rhythm} & Network/load & no & New captures and cross-network tests; overlapping windows caveat\\
MoEcho~\cite{ding2025moecho} & Task/input/user & no & Cross-template/dataset tests; contention/context conditions\\
LLM-Fin~\cite{nazari2024llmfin} & Model & no & Unseen variants within known architecture families\\
ArchSnoop~\cite{weng2026archsnoop} & Model & no & Random-weight training to official/fine-tuned models; power modes\\
WildAuth~\cite{wang2026dialogue} & Task/input/user & no & User-disjoint cross-topic split; no verified cross-service test\\
LeakyLMs~\cite{majidi2026leakylms} & Model & no & Held-out architecture configurations; known-GPU calibration\\
StepFinder~\cite{zhu2026stepfinder} & Task/input/user & no & Task-disjoint subsets; denominator reconciliation unresolved\\
Edge EM~\cite{alhalawani2026physicaledge} & Task/input/user & no & Unseen-word test separate from profiled vocabulary recovery\\
\end{longtable}


\subsection{Positive reuse and changed decisions}
Early Bird's prefix record and CachePrune's direct-recovery record are an existing descriptive match. Both can inform a cache-reuse threat shortlist; they do not support a numerical ranking across oracle and serving measurements. CachePrune supplies a configuration-linked recovery/quality/cost package for scoped feasibility analysis. The semantic processor benchmark in Early Bird instead leaves the prefix candidate's serving cost open. This produces a positive retention decision and a specific next measurement. The main decision table also distinguishes read-cover state effects, final-emitter captures and workflow utility, extending the established WF evaluation discipline to LLM/agent execution objects.
